\chapter{Исправленные опечатки}

{\small
\begin{longtable}{@{}r p{0.43\textwidth} p{0.43\textwidth}@{}}
\toprule
с. & в оригинале & исправлено\\ \midrule
\endfirsthead
\toprule
с. & в оригинале & исправлено\\ \midrule
\endhead
\bottomrule
\endlastfoot
3 & <<\dots в которых \texttt{\}}на находит приложение>> & <<\dots она находит\dots>>\\
7 & <<а в последствии>> & <<а впоследствии>>\\
11 & <<(обобщенном )>> & <<(обобщенном)>>\\
13 & <<прообраз множества $A$ при отображении $|$>> & <<\dots при отображении $f$>>\\
15 & Пример: $(A\cap\ov B\cap C)\cap(A\cap C)\equiv\dots$ & $(A\cap\ov B\cap C)\cup(A\cap C)\equiv\dots$ (иначе не проходит вынесение $A\cap C$ по $5^\circ$)\\
16 & в выкладке упр.~2: <<$([A\cap A)\cap(C\cap\ov C)]$>> & <<$[(A\cap A)\cap(C\cap\ov C)]$>>\\
19 & доказательство (1.6): сомножитель $(B\cup C\cup C)$ на четвертом месте & $(B\cup C\cup B)$\\
21 & нет номера у упражнения <<Будем предполагать\dots>> & <<4.>>\\
24 & $f[\varphi_{\alpha_1},\dots,\varphi_{\alpha_n}]$ & $f[\varphi_{\alpha_1},\dots,\varphi_{\alpha_m}]$\\
24 & <<Заменим в формуле (1.12) функции $\varphi_i$ на $\psi_i$>> & <<Заменим в этой формуле, согласно (1.12), \dots>>\\
25 & <<Рассмотрим множество $T_a\subset\M$\dots>> (в конце абзаца то же множество названо $\M_a$) & <<Рассмотрим множество $\M_a\subset\M$\dots>>\\
27 & $\varphi_0=\dfrac{x_1+x_2}{1+c^{-1}x_1x_2}$ & $\dfrac{x_1+x_2}{1+c^{-2}x_1x_2}$ (релятивистский закон сложения скоростей)\\
\end{longtable}}

Мелкие дефекты набора скана (сбитые индексы, разорванные черты
дополнения, точки вместо запятых) исправлены без отдельных отметок.
